\chapter{Software para modelización, análisis, procesamiento y síntesis de señales}

\section{Introducción a Matlab / Octave}

Matlab y Octave son intérpretes que nos permiten realizar cálculos numéricos, incluso con números complejos, vectores y matrices, y realizar gráficos en distintas representaciones.

En este apunte, mostramos algunas de las operaciones y comandos básicos para la operación interactiva. 

\subsection{La línea de comandos}

\begin{enumerate}
	
	\item Cuando Matlab/Octave comienza a ejecutarse, aparece una ventana tipo consola (sin interfaz gráfica), y una línea marcada con el símbolo \verb->>-. Esta es la línea de comandos, donde escribimos las diferentes órdenes.
	
	\item Para probar podemos escribir:
	\begin{verbatim}
		2 + 2
		3 - 5
		3 / 1.4
		2 * pi / 4
		(1 + 2 * i) / (1 - 2 * i)
	\end{verbatim}
	Nota: $i$ es la unidad imaginaria.
	
	\item En las últimas versiones de Matlab/Octave, a medida que vamos escribiendo comandos, estos quedan registrados en un historial de comandos, que podemos revisar más tarde y recordar lo hecho.
	
	\item El comando \texttt{diary} nos permite grabar en un archivo de texto todas las órdenes, datos y resultados de una sesión. Para usar el comando \texttt{diary} escribamos en la línea de comandos:
	\begin{verbatim}
		>> diary 'archivo.txt'
	\end{verbatim}
	
	\item La idea de esta clase es generar un archivo, denominado \verb-archivo.txt-, con todas las órdenes dadas.
\end{enumerate}

\subsection{El sistema de ayuda}

\begin{enumerate}
	\item Para ver cómo se usa un comando es posible escribir en la línea de comandos:
	\begin{verbatim}
		>> help comando
	\end{verbatim}
	donde \texttt{comando} es el nombre del comando que nos interesa.
	
	\item La orden \texttt{help comando} nos da una descripción del comando y posiblemente una lista de comandos relacionados.
	\begin{verbatim}
		>> help diary
	\end{verbatim}
\end{enumerate}

\subsection{Valores escalares}

\begin{enumerate}
	\item Como hemos visto, operamos con números reales y complejos, también llamados escalares.
	\item La unidad imaginaria puede escribirse como $i$ o $j$.
	\item Probemos los siguientes comandos:
	\begin{verbatim}
		(1 + 2 * j) * (1 - 2 * i)
		45 * (56 + i) - 34 * (16 - 32 * j)
	\end{verbatim}
\end{enumerate}

\subsection{Valores matriciales o arreglos}

\begin{enumerate}
	\item Matlab/Octave trabaja también con arreglos de reales y complejos.
	\item Podemos crear un vector fila escribiendo:
	\begin{verbatim}
		[ 1 + 2 * i  1 - 3 * i ]
	\end{verbatim}
	o separando con comas:
	\begin{verbatim}
		[ 1 + 2 * i, 1 - 3 * i ]
	\end{verbatim}
	
	\item Podemos crear un vector columna escribiendo:
	\begin{verbatim}
		[ 1 + 2 * i; 1 - 3 * i ]
	\end{verbatim}
	o separando con un salto de línea:
	\begin{verbatim}
		[ 1 + 2 * i
		1 - 3 * i ]
	\end{verbatim}
\end{enumerate}

\subsection{Metáfora de Matlab / Octave}

Podemos suponer, a modo de metáfora, que disponemos de un pizarrón donde escribimos las cuentas que vamos haciendo, línea por línea.

\subsection{Variables}

\begin{enumerate}
	
	\item A diferencia de una calculadora, Matlab/Octave permite asignar nombres a los diferentes valores. Esta combinación de nombre y valor recibe el nombre de variable.
	
	\item A continuación vemos ejemplos de los distintos tipos de variables. La constante $i$ es la unidad imaginaria, y ya está definida por defecto. El símbolo \texttt{\%} se usa para indicar que aquello escrito a su derecha hasta el fin de la línea es un comentario y no será ejecutado por el intérprete.
	
	\item Pruebe ingresar lo que aparece después de \verb->>-:
	\begin{verbatim}
		>> a = 5            % escalar real
		>> b = 5 + 10 * i   % escalar complejo
		>> c = [1 2 3]      % vector fila
		>> d = [1, 2, 3]    % otra forma de escribir un vector fila
		>> e = [1; 2; 3]    % un vector columna
		>> f = [1, 2, 3; 4, 5, 6; 7, 8, 9] % una matriz 3 x 3
	\end{verbatim}
\end{enumerate}

\subsection{Continuando con la metáfora de Matlab / Octave}

\begin{enumerate}
	\item Las variables son similares a divisiones que dibujamos en nuestro pizarrón, en las cuales escribimos los valores y los nombres respectivos.
	\item El espacio ocupado por las variables en memoria se denomina espacio de trabajo (workspace).
\end{enumerate}

\subsubsection{Los comandos who y what}

\begin{enumerate}
	\item El comando \texttt{who} da una lista de las variables presentes en el espacio de trabajo actual.
	\item Pruebe ingresar:
	\begin{verbatim}
		>> help who
	\end{verbatim}
	
	\item El comando \texttt{what} da una lista de archivos de código presentes en el directorio de trabajo actual. También puede probar \texttt{ls} o \texttt{dir} para listar todos los archivos. Verifique cuál es el directorio por defecto.
	\item Pruebe ingresar:
	\begin{verbatim}
		>> help what
	\end{verbatim}
	
	\item Ejecute:
	\begin{verbatim}
		>> who
		>> what
	\end{verbatim}
\end{enumerate}

\subsection{Operaciones con números complejos}

\begin{enumerate}
	\item A continuación vemos algunas de las operaciones que pueden realizarse con números complejos. Pruebe ingresar:
	\begin{verbatim}
		>> x = 3 + 4 * i
		>> real(x)  % 3
		>> imag(x)  % 4
		>> conj(x)  % 3 - 4 * i
		>> abs(x)   % 5
		>> angle(x) % 0.9273 (en radianes)
	\end{verbatim}
\end{enumerate}

\subsection{Raíces y coeficientes de un polinomio}

\begin{enumerate}
	\item Supongamos que tenemos un polinomio general:
	\begin{equation}
		b_{n}x^{n}+b_{n-1}x^{n-1}+...+b_{0}=0.
	\end{equation}
	Para hallar las raíces del polinomio podemos colocar sus coeficientes en un vector ordenado (de mayor a menor grado) y usar la función \texttt{roots}:
	\begin{verbatim}
		p = [bn, ..., b0];
		roots(p)
	\end{verbatim}
	
	\item Supongamos que tenemos que hallar las raíces del polinomio $x^{2}+2x+3=0$. Escribamos entonces:
	\begin{verbatim}
		p = [1, 2, 3];
		roots(p)
	\end{verbatim}
	
	\item Supongamos ahora que conocemos las raíces y necesitamos hallar los coeficientes del polinomio originario. Podemos usar la función \texttt{poly}:
	\begin{verbatim}
		>> r = [-1.0, 2.7, -3.1, -3.1];
		>> poly(r)
	\end{verbatim}
	
	\item Por ejemplo, si las raíces son $1, -1+2i, -1-2i$, entonces:
	\begin{verbatim}
		>> r = [1, -1+2*i, -1-2*i];
		>> poly(r)
	\end{verbatim}
	
	\item Pruebe con el vector correspondiente al polinomio $x^4+2x^3+3x^2+4x+5$:
	\begin{verbatim}
		>> p = [1, 2, 3, 4, 5]; 
		>> r = roots(p)
		>> poly(r)
	\end{verbatim}
	
	\item ¿Cuál le parece que es el resultado de la siguiente orden?
	\begin{verbatim}
		poly(roots(p))
	\end{verbatim}
	
	\item Pruebe las siguientes órdenes para mayor información:
	\begin{verbatim}
		>> help roots
		>> help poly
	\end{verbatim}
\end{enumerate}

\subsection{Generación de vectores y matrices}

\begin{enumerate}
	\item Matlab permite la generación de secuencias de valores y su asignación a vectores mediante un mecanismo sumamente eficiente basado en el operador dos puntos (\texttt{:}).
	\item Por ejemplo, supongamos que deseamos obtener un vector fila con una lista de números que comienza en $1$, termina en $2$ con incrementos de $0.1$. Escribiendo:
	\begin{verbatim}
		>> t = 1:0.1:2
	\end{verbatim}
	se obtiene el vector deseado. Observe que $t$ está formado por una fila de once columnas ($1 \times 11$).
	
	\item Si deseamos generar $t$ como un vector columna, podemos transponerlo agregando un apóstrofe (\texttt{'}):
	\begin{verbatim}
		>> t = (1:0.1:2)'
	\end{verbatim}
	
	\item Practique los siguientes puntos:
	\begin{enumerate}
		\item Genere un vector que contenga valores entre $-2$ y $2$ espaciados cada $0.01$.
		\item Para generar matrices llenas de ceros o unos podemos usar las funciones \texttt{zeros} y \texttt{ones}, respectivamente. Encuentre cómo se usan estas funciones con la ayuda (\texttt{help}) y haga un ejemplo con cada una.
	\end{enumerate}
\end{enumerate}

\subsection{Operaciones sobre vectores y matrices, primera parte}

\begin{enumerate}
	\item Matlab permite realizar operaciones con matrices siempre que sus dimensiones sean compatibles. Primero veremos cuál es la función que nos permite averiguar las dimensiones de una variable.
	\item Pruebe los siguientes comandos:
	\begin{verbatim}
		>> t = 1:0.1:2;
		>> size(t)
		>> size(t')
	\end{verbatim}
	
	\item Las operaciones algebraicas de suma, resta y multiplicación por escalares son muy simples. Pruebe:
	\begin{verbatim}
		>> t + 2
		>> t - 2
		>> t * 2
	\end{verbatim}
	
	\item Estas operaciones se realizan sumando, restando o multiplicando cada elemento individual del vector por el escalar. De la misma forma se aplican funciones trigonométricas elemento a elemento:
	\begin{verbatim}
		>> cos(t)
		>> sin(2 * t - 1)
	\end{verbatim}
	
	\item Supongamos ahora que queremos elevar al cuadrado cada elemento de un vector. Pruebe con los comandos:
	\begin{verbatim}
		>> t .* t
		>> t .^ 2
	\end{verbatim}
	
	\item Para realizar el producto matricial interno y externo de dos vectores escriba:
	\begin{verbatim}
		>> t' * t
		>> t * t'
	\end{verbatim}
	Indique la diferencia estructural entre los resultados de ambas operaciones.
\end{enumerate}

\subsection{Operaciones sobre vectores y matrices, segunda parte}

Es posible realizar operaciones vectoriales forzando el cálculo "elemento a elemento" (point-wise) agregando un punto (\texttt{.}) antes del operador matemático.
\begin{enumerate}
	\item Definamos:
	\begin{verbatim}
		x = [-1 2]
		y = [3 4]
	\end{verbatim}
	
	\item Si calculamos \texttt{z = x + y} resulta \texttt{z = [(-1 + 3) , (2 + 4)]} = \texttt{[2, 6]}.
	\item Si calculamos \texttt{z = x * y} obtendremos un mensaje de error de álgebra lineal, diciendo que las matrices no tienen dimensiones internas compatibles para ser multiplicadas.
	\item ¿Cómo podemos obtener el producto elemento a elemento $[-1 \cdot 3, 2 \cdot 4]$? Empleando el operador \texttt{.*}:
	\begin{verbatim}
		x .* y
	\end{verbatim}
	y obtendrá:
	\begin{verbatim}
		[-3 8]
	\end{verbatim}
	
	\item Lo mismo se debe hacer con la división:
	\begin{verbatim}
		x ./ y
	\end{verbatim}
	y la exponenciación:
	\begin{verbatim}
		x .^ y
	\end{verbatim}
	o
	\begin{verbatim}
		x .^ 2
	\end{verbatim}
\end{enumerate}

\subsection{Operaciones sobre vectores y matrices, tercera parte}

\begin{enumerate}
	
	\item Ingresemos un vector:
	\begin{verbatim}
		v = [1, 2, 3, 4, 5, 6, 7]
	\end{verbatim}
	
	\item Podemos acceder a los valores almacenados empleando la notación \verb-v(i)-, donde $i$ es un número entero (índice) mayor o igual a $1$. Pruebe:
	\begin{verbatim}
		v(1)
	\end{verbatim}
	
	\item También podemos extraer una porción (slice) del vector $v$ escribiendo \verb-v(i:j)-, donde $i \leq j$. Pruebe:
	\begin{verbatim}
		v(2:7)
		v(1:6)
		v(2:7) - v(1:6)  % Útil para calcular diferencias finitas
	\end{verbatim}
	
	\item Observe que la palabra reservada \texttt{end} permite acceder al último elemento dinámicamente:
	\begin{verbatim}
		v(2:end) - v(1:end-1)
	\end{verbatim}
	
	\item Matlab/Octave tiene una extensa librería de funciones de álgebra lineal aplicables a matrices. A continuación mostramos una pequeña tabla representativa:
	\begin{center}
		\begin{tabular}{ll}
			\textbf{Comando} & \textbf{Resultado} \\ \hline
			\texttt{det} & Determinante de una matriz cuadrada \\
			\texttt{eig} & Autovalores y autovectores \\
			\texttt{inv} & Inversa matricial \\
			\texttt{svd} & Descomposición en valores singulares (SVD) \\
			\texttt{cond} & Número de condición de una matriz \\
		\end{tabular}
	\end{center}
	
	\item Supongamos un sistema de 2 ecuaciones lineales con dos incógnitas representable matricialmente como $Ax=b$, donde:
	\begin{verbatim}
		A = [1, 2; 2, 1];
		b = [1; 2];
	\end{verbatim}
	
	\begin{enumerate}
		\item Matemáticamente, este sistema se resuelve calculando $x = A^{-1}b$. En código podríamos escribir:
		\begin{verbatim}
			x = inv(A) * b;
		\end{verbatim}
		\textbf{Nota importante:} Aunque usar \texttt{inv()} es correcto teóricamente, numéricamente es ineficiente y propenso a errores de precisión en matrices grandes. El método recomendado y optimizado en Matlab/Octave es usar el operador de "división izquierda" (\texttt{\textbackslash}):
		\begin{verbatim}
			x = A \ b;
		\end{verbatim}
		
		\item ¿Cómo puede verificar que el resultado es correcto? Calculando el vector de error (residuo):
		\begin{verbatim}
			e = A * x - b;
			e' * e   % Suma de los errores al cuadrado (debe ser muy cercana a 0)
		\end{verbatim}
		
		\item Resuelva los siguientes sistemas de ecuaciones y calcule los errores correspondientes:
		\begin{eqnarray}
			A_{1} = \begin{bmatrix} 1 & 1 \\ -1 & 1 \end{bmatrix}, & b_{1} = \begin{bmatrix} 0 \\ 1 \end{bmatrix}, \\
			A_{2} = \begin{bmatrix} 100.0 & 100.0 \\ 100.0 & 100.0 + 5 \cdot eps \end{bmatrix}, & b_{2} = \begin{bmatrix} 0 \\ 1 \end{bmatrix}.
		\end{eqnarray}
		
		\item ¿Para qué sirve la constante integrada \texttt{eps}? (Pista: es la distancia desde $1.0$ hasta el próximo número de punto flotante más grande en precisión doble).
		
		\item Busque en la ayuda el listado completo de operaciones matriciales (\texttt{help matfun}).
	\end{enumerate}
\end{enumerate}

\subsection{Graficación de funciones 2D}

\begin{enumerate}
	\item Para graficar funciones existen distintos comandos. Escriba:
	\begin{verbatim}
		>> t = 0:0.1:2*pi;
		>> y = cos(2 * t - 1);
		>> figure(1); plot(t, y);
		>> figure(2); stem(t, y);
	\end{verbatim}
	
	\item El comando \texttt{figure} crea o selecciona ventanas gráficas independientes. \texttt{plot} dibuja una señal de tiempo continuo uniendo los puntos con líneas, mientras que \texttt{stem} grafica la secuencia de valores discretos explícitos (ideal para señales muestreadas).
	
	\item Estos comandos admiten un tercer parámetro opcional para configurar el estilo (\verb-plot(t, y, 'formato')-). El formato es una cadena de texto (ej: \texttt{'r*: '}) que combina:
	\begin{itemize}
		\item \textbf{Colores:} \texttt{y} (amarillo), \texttt{m} (magenta), \texttt{c} (cyan), \texttt{r} (rojo), \texttt{g} (verde), \texttt{b} (azul), \texttt{w} (blanco), \texttt{k} (negro).
		\item \textbf{Marcadores:} \texttt{.} (punto), \texttt{o} (círculo), \texttt{x} (equis), \texttt{*} (asterisco), \texttt{s} (cuadrado).
		\item \textbf{Líneas:} \texttt{-} (sólida), \texttt{:} (punteada), \texttt{-.} (raya y punto), \texttt{--} (discontinua).
	\end{itemize}
	
	\item Por ejemplo:
	\begin{verbatim}
		>> plot(t, y, 'c+:')
	\end{verbatim}
	dibuja una línea punteada celeste uniendo las cruces ubicadas en cada dato.
	
	\item Observe los comandos complementarios para enriquecer el gráfico:
	\begin{verbatim}
		>> t = 0:0.1:2*pi;
		>> clf; plot(t, cos(t), 'm:o');
		>> grid on;
		>> xlabel('Tiempo (s)');
		>> ylabel('Amplitud (V)');
		>> title('Función periódica');
	\end{verbatim}
	
	\item El comando \texttt{subplot} permite colocar múltiples gráficos en la misma ventana mediante una cuadrícula (filas, columnas, índice activo). Pruebe:
	\begin{verbatim}
		>> clf;
		>> subplot(3, 2, 1); plot(t, cos(t));
		>> subplot(3, 2, 2); stem(t, cos(t));
		>> subplot(3, 2, 3); plot(t, sin(t));
		>> subplot(3, 2, 4); stem(t, sin(t));
		>> subplot(3, 2, 5); plot(t, tan(t));
		>> subplot(3, 2, 6); stem(t, tan(t));
	\end{verbatim}
	
\end{enumerate}

\subsection{Cómo guardar gráficos en archivos (Matlab y Octave)}

\begin{enumerate}
	
	\item Supongamos que hemos realizado un gráfico y queremos guardarlo como imagen para incluirlo en nuestro documento LaTeX. El formato moderno más práctico es PNG (o PDF vectorial).
	
	\item Una vez dibujada la figura, pruebe ejecutar:
	\begin{verbatim}
		>> print('mi_grafico.png', '-dpng');
	\end{verbatim}
	Esto generará el archivo visual en el directorio de trabajo actual. 
\end{enumerate}

\subsection{Graficación de funciones en tres dimensiones}

\begin{enumerate}
	\item Es posible graficar curvas paramétricas y superficies tridimensionales usando herramientas como \texttt{plot3}, \texttt{mesh} o \texttt{surf}.
	
	\item Busque la ayuda ejecutando \texttt{help plot3} o \texttt{help surf}.
	
	\item Desarrolle un script breve para dibujar una superficie en 3D (ej: usando la función \texttt{peaks} integrada) y guárdelo en un archivo de imagen.
\end{enumerate}


\subsection{Escritura y salida de texto formateado}

\begin{enumerate}
	\item Con el comando \texttt{fprintf} podemos escribir texto y valores en la pantalla de forma estructurada.
	
	\item Compare el resultado de:
	\begin{verbatim}
		>> fprintf('hola')
	\end{verbatim}
	con:
	\begin{verbatim}
		>> fprintf('hola\n')
	\end{verbatim}
	El carácter especial \verb-\n- indica el salto de línea.
	
	\item Puede imprimir valores de variables usando comodines (\texttt{\%d} para enteros, \texttt{\%f} para flotantes, \texttt{\%s} para cadenas):
	\begin{verbatim}
		>> fprintf('El valor de pi es aproximadamente %f\n', pi);
	\end{verbatim}
\end{enumerate}

\subsection{Estructuras de Control: Iteración}

\begin{enumerate}
	\item Supongamos que deseamos armar una tabla de la función $y = 2x - 1$ para todos los valores enteros de $x$ entre $1$ y $8$.
	
	\item Podemos usar un bucle \texttt{for}:
	\begin{verbatim}
		for x = 1:8
		fprintf('x = %d; y = %d\n', x, 2 * x - 1);
		end
	\end{verbatim}
	
	\item Alternativamente, usando un bucle \texttt{while}:
	\begin{verbatim}
		x = 1;
		while (x <= 8)
		fprintf('x = %d; y = %d\n', x, 2 * x - 1);
		x = x + 1; 
		end
	\end{verbatim}
\end{enumerate}

\subsection{Toma de decisiones}

\begin{enumerate}
	
	\item La función \texttt{mod(x, y)} calcula el resto de la división entera.
	\item Podemos condicionar la ejecución de código empleando \texttt{if}, \texttt{elseif}, y \texttt{else}:
	\begin{verbatim}
		for x = 1:15
		if (mod(x, 2) == 0)
		fprintf('%d es par\n', x);
		elseif (mod(x, 3) == 0)
		fprintf('%d es impar y divisible por 3\n', x);
		else
		fprintf('%d es impar\n', x);
		end
		end
	\end{verbatim}
\end{enumerate}

\section{Ejemplo 1: Integración gráfica}

En esta sección repasaremos los conceptos escribiendo un pequeño script para:
\begin{enumerate}
	\item Generar un vector eje de abscisas de 51 números.
	\item Generar datos aleatorios asociados a ese eje.
	\item Graficar con etiquetas correspondientes al título y los ejes.
\end{enumerate}

Para graficar los números generados, escriba el siguiente código y ejecútelo:
\begin{verbatim}
	x = 1:1:51;            % Los datos del eje temporal discreto
	y = rand(size(x));     % Señal de ruido aleatorio uniforme entre 0 y 1
	plot(x, y, 'b-*');
	title('Señal estocástica generada por computadora');
	xlabel('Índice de muestra (n)');
	ylabel('Amplitud aleatoria');
	grid on;
\end{verbatim}

\section{Ejemplo 2: Interacción con archivos de texto (I/O)}

En esta sección veremos cómo escribir y leer datos estructurados.
\begin{enumerate}
	\item Usamos la función \texttt{fopen} para abrir o crear el archivo, obteniendo un Identificador (\texttt{fid}):
	\begin{verbatim}
		fid = fopen('datos.txt', 'w'); % 'w' para escritura, 'r' lectura
	\end{verbatim}
	
	\item A continuación, el código de escritura (empleando bucles dobles anidados) y lectura vectorizada:
	\begin{verbatim}
		% -- ESCRITURA --
		esc = fopen('matriz_datos.txt', 'w');
		
		for fila = 1:3
		for col = 1:10
		fprintf(esc, '%d ', fila + col); % Espacio separador
		end
		fprintf(esc, '\n'); % Salto al terminar la fila
		end
		fclose(esc); % SIEMPRE liberar el recurso
		
		% -- LECTURA --
		lec = fopen('matriz_datos.txt', 'r');
		% Lee los datos respetando la estructura 3x10 definida
		[A, cantidad_leida] = fscanf(lec, '%d', [3, 10]);
		fclose(lec);
	\end{verbatim}
	
	\item La variable \texttt{cantidad\_leida} resultará ser $30$, verificando que no hubo errores en el proceso.
\end{enumerate}

\section{Ejemplo 3: Vectorización eficiente y protección matemática}

En el siguiente algoritmo procesaremos datos evitando el uso ineficiente de ciclos \texttt{for} y protegiéndonos de cálculos indefinidos (como divisiones por cero).

\begin{enumerate}
	\item La indexación lógica \verb-(abs(vector) < eps)- nos permite encontrar qué elementos son peligrosamente cercanos a cero.
	\item \texttt{eps} es la constante de "precisión de máquina" en punto flotante, el menor número manejable para comparaciones con el cero.
	
	\begin{verbatim}
		% Supongamos que ya leímos una matriz A de Nx3 dimensiones
		a = A(:, 1);  % Extrae TODA la columna 1
		b = A(:, 2);
		c = A(:, 3);
		
		% Protegemos el divisor 'c' de valer exactamente 0
		ud = c + (abs(c) < eps) * eps; 
		vd = (a - b);
		vd = vd + (abs(vd) < eps) * eps;
		
		% Cálculos vectorizados: División punto a punto simultánea
		u = (a + b) ./ ud;  
		v = c ./ vd;        
	\end{verbatim}
\end{enumerate}